\scanpage{147}
Considérons le sous-groupe de \deleted{\unsure{ce}} $D_{\mathbb R\times\mathbb U}$ formé des $f_{a,\theta}$ qui \deleted{ont} sur $(\partial X)_{-1}$ et $(\partial X)_{+1}$ sont des \deleted{\unsure{auto}} op. d'ordre fini, ce qui s'exprime par les conditions
\[
\theta\exp-2i\pi a,\quad\theta\exp2i\pi a\quad\in\deleted{\unsure{\infty}}\mu_\infty(\mathbb C)
\]
\margin{(ens. des él. d'ordre fini de $\mathbb U$)}
soit encore
\[
(28)\qquad\theta\in\mu_\infty(\mathbb C),\quad a\in\mathbb Q.
\]
Ce sous-groupe est stable par l'opération de $\{\pm1\}\times\{1,\sigma\}$, opérant de la façon explicitée plus haut -- donc on trouve un ss-groupe de (24), formé \inserted{justement} par les opérations \inserted{(\unsure{dudit gr|du total gr}} qui sur $\partial X$ sont d'ordre fini. Ce sous-groupe \deleted{est} s'identifie :
\[
(29)\qquad G=\deleted{\unsure{\{\pm1\}\ltimes}}(\mu\times\mu).(\mathbb Q\times\mu_\infty)\qquad1/2\text{ direct}
\]
où $\mu=\{\pm1\}$, $\mu_\infty=\mu_\infty(\mathbb C)\simeq\mathbb Q/\mathbb Z$
\[
\exp2i\pi q\longleftarrow\!\shortmid\ q\bmod\mathbb Z
\]
et où $\mu\times\mu$ opère sur $\mathbb Q\times\mu_\infty$ par symétrie sur les deux facteurs. On identifie bien sûr $\sigma$ à $(-1,-1)\in\mu\times\mu$. On pose
\[
(30)\qquad\deleted{\unsure{G}}\widetilde G=(\mu\times\mu).(\mathbb Q\times\mathbb Q)\qquad1/2\text{ direct}
\]
avec une suite exacte
\diagram{\begin{tikzcd}
(31)\quad 1 \arrow[r,dash] & \mathbb Z \arrow[r] & \deleted{\unsure{G}}\widetilde G \arrow[r] & G \arrow[r] & 1\\
& n \arrow[r,mapsto] & \unsure{\{f_{1,n}\}|f_{0,n}} & & \\
& & (\varepsilon,\varepsilon').\unsure{\{f_{a,b}\}|f_{a,b}} \arrow[r,mapsto] & (\varepsilon,\varepsilon').\unsure{\{g_{a,\exp2i\pi b}\}|g_{a,\exp2i\pi b}} &
\end{tikzcd}}
\margin{147}
